\begin{table}[H]
\centering
\caption{Relative performance comparison by MEAN_SPL using the Friedman test with post-hoc Nemenyi analysis.}
\label{tab:stats_mean_spl}
\begin{threeparttable}
\begin{tabular}{lccccc}
\toprule
Model & Mean MEAN_SPL & Std. & Avg. rank & Sig. wins & Sig. losses \\
\midrule
SetTransformer\_Q & 0.2106 & 0.0644 & 3.207 & 6 & 0 \\
M1\_AggHistFutureSummary & 0.2179 & 0.0686 & 3.545 & 6 & 0 \\
DeepSets\_Q & 0.2236 & 0.0829 & 3.793 & 6 & 0 \\
M3\_FullSkuTemporalCNN\_Q & 0.2203 & 0.0788 & 3.851 & 6 & 0 \\
M0\_AggHistOnly & 0.2181 & 0.0684 & 3.926 & 6 & 0 \\
BottomUpGlobalHistGB & 0.2706 & 0.1294 & 6.595 & 3 & 5 \\
AggregateHistGB & 0.2641 & 0.1019 & 7.140 & 2 & 5 \\
ChildSummaryHistGB & 0.2632 & 0.1009 & 7.207 & 2 & 5 \\
AggregateElasticNet & 0.3814 & 0.2908 & 8.223 & 1 & 6 \\
SeasonalNaive & 0.3412 & 0.1826 & 8.992 & 0 & 8 \\
ChildSummaryElasticNet & 0.4211 & 0.3281 & 9.521 & 0 & 9 \\
\bottomrule
\end{tabular}
\begin{tablenotes}
\footnotesize \item The Friedman test uses 121 aligned series-level blocks (task,agg\_id). The test statistic is $\chi^2=627.1736$ with $p=2.641e-128$. Average ranks are computed with lower MEAN_SPL indicating better performance. Nemenyi significant wins/losses are counted at $\alpha=0.10$ using critical difference $CD=1.2697$.
\end{tablenotes}
\end{threeparttable}
\end{table}
